\capitulo{7}{Conclusiones y Líneas de trabajo futuras}

El resultado del proyecto comprende tanto la aplicación desarrollada como
el aprendizaje adquirido al construirla, y deja abiertas distintas
posibilidades para continuar trabajando en \textit{Manualito}.

\section{Conclusiones}
\label{conclusiones:balance}

Al relacionar el trabajo realizado con los objetivos planteados, se pueden
extraer las siguientes conclusiones:

\begin{itemize}
    \item Se ha desarrollado una aplicación web que reúne la gestión de
    juegos, manuales y conversaciones con la generación de explicaciones
    y la consulta de dudas, dando forma al objetivo de ofrecer una
    herramienta de apoyo antes y durante una partida.

    \item La ejecución local de los modelos ha requerido adaptar las
    decisiones técnicas a los recursos disponibles, por lo que la
    investigación y las pruebas han formado parte del desarrollo.
    Comparar alternativas ha servido para valorar tanto sus resultados
    como las condiciones necesarias para incorporarlas a la aplicación.

    \item El proyecto ha permitido aplicar conocimientos del grado y
    profundizar en el procesamiento de documentos, la inteligencia
    artificial y el desarrollo web. Al abordar también la interfaz,
    la seguridad, las pruebas y el despliegue, se ha adquirido una visión
    más completa del trabajo necesario para publicar una aplicación.

    \item El uso supervisado de agentes ha servido para explorar el
    desarrollo agéntico en tareas concretas, desde la escritura de código
    repetitivo hasta la elaboración de pruebas. Su utilización ha requerido
    comprender y comprobar las soluciones propuestas antes de incorporarlas,
    manteniendo el criterio propio sobre la arquitectura y el diseño.

    \item La organización por iteraciones ha ayudado a revisar las
    prioridades y adaptar el trabajo a las necesidades que iban apareciendo.
    Esta experiencia refuerza la importancia de acotar el alcance y reservar
    tiempo para integrar, probar y pulir las funcionalidades desarrolladas.
\end{itemize}

\section{Líneas de trabajo futuras}
\label{conclusiones:trabajo-futuro}

El tiempo disponible ha limitado el alcance del proyecto, por lo que seguir
trabajando en \textit{Manualito} permitiría retomar algunas funcionalidades
pendientes y abordar nuevas ampliaciones como las siguientes:

\begin{itemize}
    \item \textbf{Administración y moderación de manuales.} Desarrollar
    un panel para gestionar usuarios y revisar los manuales subidos, de
    modo que puedan retirarse los archivos que no correspondan a un manual
    o cuya calidad impida leer sus páginas o empeore notablemente la calidad de las respuestas del \ac{LLM}.

    \item \textbf{Recomendación de juegos.} Durante el proyecto se exploró,
    como iniciativa paralela, un recomendador basado en \ac{ML} para ayudar
    al usuario a descubrir nuevos juegos, pero faltó tiempo para pulirlo
    e incorporarlo a la entrega. Esta línea podría
    retomarse revisando las recomendaciones obtenidas y completando los
    ajustes y las pruebas necesarios para valorar su integración.

    \item \textbf{\acs{TTS}.}\acused{TTS} La síntesis de voz permitiría
    escuchar las explicaciones y respuestas sin consultar continuamente la
    pantalla, aunque se descartó por las limitaciones del \textit{hardware}
    disponible. Retomarla requeriría estudiar una solución que genere el
    audio sin perjudicar el funcionamiento del resto de la aplicación.
\end{itemize}
